\section{变长指令编码与逐字节审计}

x86 没有统一固定长度。一个指令可以包含前缀、操作码、ModR/M、SIB、位移和
立即数，最长不超过 15 字节。当前执行模式给出默认操作数与地址宽度，前缀再
覆盖默认值。译码器必须从字节流判断当前指令在哪里结束，下一条指令才从正确
边界开始。

\begin{osgridlongtable}{L{0.15\textwidth}L{0.17\textwidth}L{0.23\textwidth}L{0.33\textwidth}}
\textbf{字段} & \textbf{是否总有} & \textbf{典型内容} & \textbf{作用}\\ \hline
\endfirsthead
\textbf{字段} & \textbf{是否总有} & \textbf{典型内容} & \textbf{作用}\\ \hline
\endhead
前缀 & 否 & 段覆盖、宽度、重复 & 修改默认译码或内存来源\\ \hline
操作码 & 是 & 一至数个字节 & 选择指令语义\\ \hline
ModR/M & 否 & mod、reg、r/m & 选择寄存器和寻址形态\\ \hline
SIB & 否 & scale、index、base & 表达比例索引地址\\ \hline
位移 & 否 & 8/16/32 位 & 加到有效地址或相对目标\\ \hline
立即数 & 否 & 指令内嵌常量 & 作为源操作数或掩码\\ \hline
\end{osgridlongtable}

变长编码来自长期兼容扩展。8086 操作码不能全部重新分配，新功能只能复用转义
字节、前缀和新增字段。它提升了代码密度，也让错误控制流危险：跳到一条指令
中间后，同一串字节可能被重新分割成另一组合法指令。

\topic{ModR/M 把寄存器与内存形式压入一个字节}
ModR/M 的高两位 mod 选择寄存器或不同位移长度，中间三位 reg 通常选择一个
寄存器或扩展操作码，低三位 r/m 选择另一个寄存器或寻址组合。16 位地址模式
中的 r/m 表对应 BX、BP、SI、DI 的历史组合，32/64 位模式则配合 SIB 表达
\(base+index\times scale+displacement\)。

例如 16 位 \code{mov ds,ax} 编码为 \code{8E D8}。\code{8E} 选择从
r/m16 装入段寄存器，\code{D8}=二进制 \code{11 011 000}：mod 为
\code{11} 表示寄存器，reg 为 \code{011} 选择 DS，r/m 为
\code{000} 选择 AX。两字节共同表达完整语义。

\topic{立即数宽度由指令形式而不是数值大小决定}
\code{mov sp,0x7000} 在 16 位模式编码为 \code{BC 00 70}。\code{BC}
已经把目标寄存器 SP 编入操作码低位，后面跟 16 位立即数。即使数值较小，
该形式仍按操作数宽度保存两个字节；小端序使低字节 \code{00} 在前。

汇编器有时可以在多个等价形式中选择较短编码，例如有符号 8 位立即数与完整
宽度立即数。源码中的数值相同不保证机器码长度相同，因此依赖固定偏移的复位
区域必须由链接断言和产物审计保护，不能按源代码行数估计大小。

\topic{操作数宽度、地址宽度和指针宽度相互独立}
操作数宽度决定一次算术或搬运处理多少位；地址宽度决定有效地址计算使用哪些
寄存器和回绕规则；指针宽度是软件类型与 ABI 的约定。16 位入口可以用前缀
执行 32 位操作，64 位模式也可以访问 8 位设备寄存器。

宽度扩大时要区分零扩展与符号扩展。无符号地址通常零扩展，有符号位移必须
符号扩展。把 \code{0xF00D} 当无符号 61453 加到 IP，与把它解释为
\(-4083\) 得到的低 16 位相同，但进入更宽计算后语义会分离。公式和类型必须
明确扩展规则。

\topic{标志寄存器连接算术与控制流}
加减和比较会更新 CF、ZF、SF、OF 等状态。CF 描述无符号进位或借位，OF
描述有符号结果越界，ZF 描述结果为零，SF 复制结果最高位。相同位模式可以在
有符号与无符号比较中产生不同结论，所以条件跳转分为 \code{ja/jb} 与
\code{jg/jl} 等两组。

\begin{osgridlongtable}{L{0.18\textwidth}L{0.23\textwidth}L{0.24\textwidth}L{0.23\textwidth}}
\textbf{场景} & \textbf{关键标志} & \textbf{条件跳转} & \textbf{解释域}\\ \hline
\endfirsthead
\textbf{场景} & \textbf{关键标志} & \textbf{条件跳转} & \textbf{解释域}\\ \hline
\endhead
相等 & ZF & \code{je/jne} & 位模式相等\\ \hline
无符号低于 & CF & \code{jb/jae} & 地址、容量、端口值\\ \hline
有符号小于 & SF 与 OF & \code{jl/jge} & 有符号整数\\ \hline
加法溢出 & OF & \code{jo/jno} & 补码范围\\ \hline
过程状态 & 项目约定 CF & \code{jc/jnc} & 汇编内部 ABI\\ \hline
\end{osgridlongtable}

当前串口过程主动用 CF 返回成功状态，这是软件约定而非 UART 自动设置。过程
中任何会修改 CF 的指令都必须安排在最终 \code{stc/clc} 之前，调用者也要在
下一条可能改标志的指令前消费结果。

\topic{CALL 与 RET 建立最小控制流栈}
近 \code{call} 把下一条 IP 压栈，再把目标位移加到 IP；\code{ret} 从
SS:SP 弹出返回 IP。16 位压栈先让 SP 减 2，再向新位置写入两个字节。若
SP 初始为 \code{0x7000}，首次返回地址占用 \code{0x6FFE..0x6FFF}。

嵌套调用形成一串返回地址。递归或过深调用会向低地址持续增长，最终覆盖其他
数据。硬件不会知道某段 RAM 被软件称为“栈”，守护页和边界检查要等分页与
内核设施建立后才能提供。

\topic{字符串指令包含隐式操作数}
\code{lodsb} 隐式读取 DS:SI 到 AL，并根据 DF 增减 SI；\code{stosb}
隐式写入 ES:DI；\code{rep movsb} 还隐式使用计数寄存器。它们代码紧凑，
却把状态藏在寄存器约定里。入口必须清 DF，调用约定必须说明哪些隐式寄存器
被破坏。

\code{cs lodsb} 的段覆盖前缀只影响本次源内存访问。编码中的前缀不是单独
执行的指令，CPU 把它和后续操作码作为一个整体译码。如果控制流落到操作码
而跳过前缀，内存来源就会变回 DS。

\topic{汇编过程也需要清晰 ABI}
每个过程都要记录输入寄存器、输出寄存器、破坏集合、标志返回、栈使用和可重入
条件。寄存器数量有限不意味着可以隐式共享所有状态。下面的最小约定适合当前
单线程固件：

\begin{osgridlongtable}{L{0.19\textwidth}L{0.31\textwidth}L{0.38\textwidth}}
\textbf{过程} & \textbf{输入与输出} & \textbf{主要破坏状态}\\ \hline
\endfirsthead
\textbf{过程} & \textbf{输入与输出} & \textbf{主要破坏状态}\\ \hline
\endhead
UART 初始化 & 无输入；CF 返回结果 & AX、DX、计数寄存器\\ \hline
发送字节 & AL 为字节；CF 返回结果 & DX、轮询计数、标志\\ \hline
发送字符串 & CS:SI 为零结尾文本 & AX、SI、DX、计数、标志\\ \hline
\end{osgridlongtable}

后续汇编调用 C++ 时要切换到平台 ABI，保存被调用者保存寄存器并保证栈对齐。
内部小过程可以使用更简单约定，但边界必须有唯一书面定义。

\topic{从源码到字节的逐层审计}
审计先用 \code{nasm -l} 或反汇编得到地址与字节，再把关键指令按字段解码。
复位入口至少核对 \code{cli}、\code{cld}、段初始化、栈设置和首个调用；
复位向量核对操作码、位移端序和目标；端口访问核对 DX 与 AL/AX 宽度。

单步调试观察的是动态路径，反汇编覆盖静态产物。未执行分支、填充区域和错误
镜像仍要靠静态审计与故障测试。两类证据结合后，才能区分“字节从未生成正确”
和“正确字节在某状态下走错路径”。
